\section{Monsters}
\label{sec:cses-monsters}

\problemstatement{On a grid with your start \texttt{A}, monsters
\texttt{M}, walls \texttt{\#}, and floor \texttt{.}, you and all monsters
move one step per turn (four directions). Escape by reaching any boundary
cell before a monster occupies it. Output \texttt{YES} and a path, or
\texttt{NO}.}

\begin{patternbox}
This is a race modelled with \textbf{multi-source BFS}. First compute, for
every cell, the earliest time a monster can arrive (BFS from all monsters
at once). Then BFS from \texttt{A}, only allowed to enter a cell
\emph{strictly before} any monster could -- reaching the boundary means
escape.
\end{patternbox}

\subsection*{Two BFS passes: monsters first, then you}

Run a BFS seeded with \emph{all} monster cells simultaneously; this labels
each cell with the minimum monster arrival time. Then BFS from \texttt{A}
with your own time counter, stepping into a neighbour only if your arrival
time there is strictly less than the monster time (so a monster cannot be
there yet). Track parents to print the path. If your BFS reaches any
boundary cell, you escape; otherwise output \texttt{NO}.

\begin{ideabox}[Multi-Source BFS Gives Simultaneous Spread]
Seeding a BFS queue with every monster at once yields, in one pass, the
soonest a monster reaches each cell -- as if all monsters spread in
parallel. Comparing your arrival time against that field decides which
cells are safe to enter.
\end{ideabox}

\subsection*{Algorithm}

\begin{enumerate}
  \item Multi-source BFS from all \texttt{M} cells to get
        $\textit{monsterTime}[\,]$.
  \item BFS from \texttt{A} tracking $\textit{myTime}$ and parents;
        enter a cell only if $\textit{myTime}<\textit{monsterTime}$.
  \item If a boundary cell is reached, reconstruct and print the path
        (\texttt{YES}); else \texttt{NO}.
\end{enumerate}

\subsection*{Dry run}

\dryrun{\texttt{A} near the centre with one distant \texttt{M}.}

\begin{itemize}
  \item Monster BFS marks cells near \texttt{M} with small times, cells
        near \texttt{A} with large times.
  \item Your BFS heads to the nearest boundary, reaching it before the
        monster's time there -- escape with a printed path.
\end{itemize}

\subsection*{C++ implementation}

\begin{lstlisting}
#include <bits/stdc++.h>
using namespace std;

int main() {
    int n, m;
    cin >> n >> m;
    vector<string> g(n);
    for (auto& row : g) cin >> row;

    const int INF = 1e9;
    vector<vector<int>> mt(n, vector<int>(m, INF));   // monster arrival time
    queue<pair<int,int>> q;
    int sr = 0, sc = 0;
    for (int i = 0; i < n; i++)
        for (int j = 0; j < m; j++) {
            if (g[i][j] == 'M') { mt[i][j] = 0; q.push({i, j}); }
            if (g[i][j] == 'A') { sr = i; sc = j; }
        }
    int dr[] = {-1,1,0,0}, dc[] = {0,0,-1,1};
    char mvc[] = {'U','D','L','R'};

    while (!q.empty()) {                              // multi-source BFS
        auto [r, c] = q.front(); q.pop();
        for (int k = 0; k < 4; k++) {
            int nr = r + dr[k], nc = c + dc[k];
            if (nr < 0 || nr >= n || nc < 0 || nc >= m) continue;
            if (g[nr][nc] == '#' || mt[nr][nc] != INF) continue;
            mt[nr][nc] = mt[r][c] + 1;
            q.push({nr, nc});
        }
    }

    vector<vector<int>> myt(n, vector<int>(m, INF));
    vector<vector<char>> mv(n, vector<char>(m, 0));
    myt[sr][sc] = 0;
    q.push({sr, sc});
    int er = -1, ec = -1;
    while (!q.empty()) {
        auto [r, c] = q.front(); q.pop();
        if (r == 0 || r == n - 1 || c == 0 || c == m - 1) { er = r; ec = c; break; }
        for (int k = 0; k < 4; k++) {
            int nr = r + dr[k], nc = c + dc[k];
            if (nr < 0 || nr >= n || nc < 0 || nc >= m) continue;
            if (g[nr][nc] == '#' || myt[nr][nc] != INF) continue;
            if (myt[r][c] + 1 >= mt[nr][nc]) continue;    // monster reaches first
            myt[nr][nc] = myt[r][c] + 1;
            mv[nr][nc] = mvc[k];
            q.push({nr, nc});
        }
    }

    if (er == -1) { cout << "NO\n"; return 0; }
    string path;
    int r = er, c = ec;
    while (!(r == sr && c == sc)) {
        char ch = mv[r][c]; path += ch;
        if (ch == 'U') r++; else if (ch == 'D') r--;
        else if (ch == 'L') c++; else c--;
    }
    reverse(path.begin(), path.end());
    cout << "YES\n" << path.size() << "\n" << path << "\n";
    return 0;
}
\end{lstlisting}

\begin{complexitybox}
\textbf{Time Complexity.} $O(nm)$ -- two BFS passes over the grid.
\textbf{Space Complexity.} $O(nm)$ for the time and move grids.
\end{complexitybox}

\begin{takeawaybox}
Races against spreading hazards are multi-source BFS (the hazard field)
plus a guarded single-source BFS (your moves), entering a cell only when
you strictly beat the hazard. This two-field comparison is the template for
fire/flood-escape problems.
\end{takeawaybox}
